\section{Piecewise linear certificates in dimension one}
\label{sec:piecewise-linear}

Consider the following setup:

\begin{setup}
Let:
\begin{enumerate}
    \item $dX_t = f(X_t) dt + g(X_t) dW_t$ be a one-dimensional SDE
    \item $V(x)$ be a piecewise linear function such that $V(x) = a_i x + b_i$, where $x \in \left( x_i, x_{i+1} \right)$.
    \item $V(x)$ is continuous at breakpoints: $V(x_i^+)=V(x_i^-) = V(x_i)$
    \item as a consequence, $V'_{-}(x_i) = a_{i-1}$, $V'_{+}(x_i) = a_i$
\end{enumerate}

\end{setup}

Now we will consider what conditions on $V$ are sufficient to guarantee that $V(X_t)$ is a supermartingale.

\begin{hypothesis}
Does concavity of $V(x)$ (along with an interior / drift term condition) guarantee that $V(X_t)$ is a supermartingale?
\end{hypothesis}

To answer this, let's apply the Itô-Tanaka formula to $V(X_t)$:
\begin{align}
    V(X_t) = V(X_0) + &\int_0^t V'(X_s) f(X_s) ds \\
        + &\int_0^t V'(X_s) g(X_s) dW_s \\
        + \frac{1}{2} &\sum_{i=1}^{n-1} \left( V'_{+}(x_i) - V'_{-}(x_i) \right) L_t^{x_i}(X) \\
\end{align}

We will focus on these terms: 

\begin{enumerate}
  \item $D_t = \int_0^t V'(X_s) f(X_s) ds$, the drift term
  \item $M_t = \int_0^t V'(X_s) g(X_s) dW_s$, the diffusion term (stochastic integral)
  \item $K_t = \frac{1}{2} \sum_{i=1}^{n-1} \left( V'_{+}(x_i) - V'_{-}(x_i) \right) L_t^{x_i}(X)$, the kink term
\end{enumerate}

we need to prove that $D_t + M_t + K_t$ is a supermartingale.

Some natural assumptions, which are sufficient to guarantee that $V(X_t)$ is a supermartingale, are the following:
\begin{enumerate}
  \item $f(x) V'(x) \leq 0$ for all $x \in (x_i, x_{i+1})$. This is because of the following:
  
  \begin{fact}
  \label{fact:nonpositive-drift-supermartingale}
  If $f(X) V'(X) \leq 0$, then the drift term $D_t$ is a supermartingale. \\
  \linebreak
  \textbf{Proof:} We want to prove $\mathbb{E}[D_t | \mathcal{F}_s] \leq D_s$ for $s \leq t$.
  \begin{align}
    \mathbb{E}[D_t | \mathcal{F}_s] &= \mathbb{E}\left[\int_0^t V'(X_u) f(X_u) du \Big| \mathcal{F}_s\right] \\
    &= \int_0^s V'(X_u) f(X_u) du + \mathbb{E}\left[\int_s^t V'(X_u) f(X_u) du \Big| \mathcal{F}_s\right] \\
    &\leq D_s + 0 = D_s
  \end{align}
  \end{fact}

  \item $M_t$ is a martingale. Assuming $V(x)$ has finitely many pieces
  (as is the case for a neural network with ReLU activation functions), it is enough
  to impose a square-integrability condition on $g(x)$\footnote{For example,
  $\mathbb{E}\!\left[\int_0^t g(X_s)^2\,ds\right]<\infty$.}
  
  \item $V'_{+}(x_i) - V'_{-}(x_i) \leq 0 \iff a_{i-1} \geq a_i$, i.e. the \textbf{concavity condition}. This is because of the following:
  \begin{fact}
    If $V'_{+}(x_i) - V'_{-}(x_i) \leq 0$ and $L_t^{x_i}(X) \geq 0$, then the kink term $K_t$ is a supermartingale. \\
    \linebreak
    \textbf{Proof:} We want to prove $\mathbb{E}[K_t | \mathcal{F}_s] \leq K_s$ for $s \leq t$.
    \begin{align}
      \mathbb{E}[K_t | \mathcal{F}_s] &= \mathbb{E}\left[\frac{1}{2} \sum_{i=1}^{n-1} \left( V'_{+}(x_i) - V'_{
-}(x_i) \right) L_t^{x_i}(X) \Big| \mathcal{F}_s\right] \\
      &= \frac{1}{2} \sum_{i=1}^{n-1} \left( V'_{+}(x_i) - V'_{-}(x_i) \right) \mathbb{E}\left[L_t^{x_i}(X) | \mathcal{F}_s\right] \\
      &\leq \frac{1}{2} \sum_{i=1}^{n-1} \left( V'_{+}(x_i) - V'_{-}(x_i) \right) L_s^{x_i}(X) = K_s
    \end{align}
    Note that the inequality follows from the fact that $L_t^{x_i}(X)$ is non-decreasing in $t$, hence $\mathbb{E}\left[L_t^{x_i}(X) | \mathcal{F}_s\right] \geq \mathbb{E}\left[L_s^{x_i}(X) | \mathcal{F}_s\right] = L_s^{x_i}(X)$.
  \end{fact}

\end{enumerate}

We will consider what happens in different practical cases to these terms, including: 
\begin{enumerate}
  \item A simple one-dimensional SDE with a piecewise linear certificate.
  \item A pendulum with a stochastic torque, and $X_t$ measuring the angle of the pendulum.
  \item A network $b$ black and $w$ white nodes, with given transition probabilities, and $X_t$ measuring the amount of mass in the white nodes. 
\end{enumerate}
to ensure our assumptions are reasonable.

\subsection{Proof of the supermartingale property}

First, let's note the following:

\begin{fact}
  The sum of a martingale and a supermartingale is a supermartingale. \\
  \linebreak
  \textbf{Proof}: By definition, given a martingale $M_t$ and a supermartingale $S_t$, for $r < t$ we have 
  $\mathbb{E}[M_t | \mathcal{F}_r] = M_r$ and 
  $\mathbb{E}[M_t | \mathcal{F}_r] \leq S_r$. Then, for $r < t$:
  \begin{align}
    \mathbb{E}[M_t + S_t | \mathcal{F}_r] &= \mathbb{E}[M_t | \mathcal{F}_r] + \mathbb{E}[S_t | \mathcal{F}_r] \\
    &= M_r + \mathbb{E}[S_t | \mathcal{F}_r] \\
    &\leq M_r + S_r = (M + S)_r
  \end{align}
\
hence $M_t + S_t$ is a supermartingale.

\end{fact}

The drift term $D_t$ is a supermartingale by assumption 1. The kink term $K_t$ is a supermartingale by assumption 3. The stochastic integral $M_t$ is a martingale by assumption 2. Therefore, the sum of these three terms is a supermartingale (as a consequence of the previous fact), and hence $V(X_t)$ is a supermartingale.

